\documentclass[twocolumn,11pt]{article}

% ═══════════ Packages ═══════════
\usepackage[utf8]{inputenc}
\usepackage[T1]{fontenc}
\usepackage[english]{babel}
\usepackage{amsmath,amssymb}
\usepackage{graphicx}
\usepackage{booktabs}
\usepackage{hyperref}
\usepackage{xcolor}
\usepackage{float}
\usepackage{enumitem}
\usepackage{authblk}
\usepackage[margin=2cm]{geometry}

\hypersetup{
  colorlinks=true,
  linkcolor=blue!70!black,
  citecolor=blue!70!black,
  urlcolor=blue!70!black,
}

% ═══════════ Title ═══════════
\title{%
  \textbf{ADFormerHybrid: A Dual-Branch Transformer Architecture \\
  with Ensemble Voting for Alzheimer's Disease Detection \\
  from Scalp EEG Signals}
}

\author[1]{Romain Kocupyr}
\affil[1]{Independent Researcher — Neuro-Link Project}
\date{February 2026}

\begin{document}

\maketitle

% ═══════════ Abstract ═══════════
\begin{abstract}
We present \textbf{ADFormerHybrid}, a novel dual-branch Transformer architecture for automated Alzheimer's disease (AD) detection and staging from standard 19-channel scalp electroencephalogram (EEG) recordings. The model combines a \textit{patch-based Transformer encoder} operating on raw EEG temporal patches with a \textit{feature encoder} processing 267 handcrafted spectral, entropy, connectivity, and asymmetry biomarkers. Cross-modal fusion is achieved through additive token injection before the Transformer layers, followed by a concatenation-based classification head. We train five independent model checkpoints and aggregate predictions via soft-voting ensemble, achieving \textbf{99.95\% screening accuracy} and \textbf{97.79\% staging accuracy} on a 65-patient clinical dataset (AD vs.\ healthy controls). The two-stage pipeline first detects AD presence, then classifies severity (mild, moderate, severe). We release Neuro-Link, an open-source clinical decision support platform integrating ADFormerHybrid with a FastAPI backend, React interface, and AI-powered report generation. Code is available at \url{https://github.com/romainsantoli-web/neuro-link}.
\end{abstract}

\noindent\textbf{Keywords:} Alzheimer's disease, EEG, Transformer, deep learning, ensemble voting, biomarkers, clinical decision support

% ═══════════ 1. Introduction ═══════════
\section{Introduction}

Alzheimer's disease (AD) affects over 55 million people worldwide, with numbers projected to triple by 2050 \cite{who2023}. Early detection remains critical, yet current diagnostic methods (PET, CSF biomarkers, MRI) are expensive, invasive, or inaccessible in primary care settings. Electroencephalography (EEG) offers a non-invasive, low-cost, and widely available alternative that captures functional brain activity in real time.

Recent advances in deep learning have demonstrated the potential of EEG-based AD detection. However, most approaches rely on either raw signal processing \cite{ieracitano2020} or handcrafted features \cite{dauwels2010}, rarely combining both modalities. Moreover, clinical deployment requires not just high accuracy but also interpretability, reliability, and integration into existing workflows.

In this work, we introduce \textbf{ADFormerHybrid}, a dual-branch architecture that simultaneously processes:
\begin{enumerate}[nosep]
  \item \textbf{Raw EEG patches} through a Transformer encoder with learnable positional embeddings,
  \item \textbf{267 handcrafted biomarkers} through a dedicated feature encoder with layer normalization.
\end{enumerate}

Our contributions are threefold:
\begin{itemize}[nosep]
  \item A novel cross-modal fusion mechanism that injects biomarker representations into the Transformer attention space;
  \item A five-model soft-voting ensemble strategy that improves staging accuracy by +1.15\%;
  \item An end-to-end open-source clinical platform (Neuro-Link) with AI-generated reports.
\end{itemize}

% ═══════════ 2. Related Work ═══════════
\section{Related Work}

\textbf{Traditional EEG biomarkers.}
Spectral power ratios (theta/alpha), coherence, and entropy measures have been extensively studied as AD biomarkers \cite{dauwels2010, jeong2004}. While interpretable, standalone features rarely achieve above 90\% accuracy on heterogeneous datasets.

\textbf{Deep learning approaches.}
CNNs \cite{ieracitano2020}, RNNs \cite{morabito2017}, and more recently Vision Transformers (ViT) \cite{dosovitskiy2020} have been applied to EEG classification. EEGNet \cite{lawhern2018} established a compact architecture for BCI tasks, but was not designed for neurodegenerative disease detection.

\textbf{Hybrid methods.}
Some works combine handcrafted features with learned representations \cite{li2022}, but typically through late fusion (concatenation at the classification head). Our approach differs by performing \textit{early fusion} at the token level before attention computation.

% ═══════════ 3. Methods ═══════════
\section{Methods}

\subsection{Data Acquisition and Preprocessing}

We use 19-channel EEG recordings from 65 patients (AD and healthy controls) sampled at 128~Hz according to the international 10-20 system. Preprocessing includes:
\begin{itemize}[nosep]
  \item Bandpass filtering: 0.5--45~Hz
  \item Notch filter: 50~Hz (power line noise)
  \item Segmentation: non-overlapping 4-second epochs ($512 \times 19$ samples)
  \item Z-score normalization per channel
\end{itemize}

\subsection{Feature Extraction (267 dimensions)}

For each segment, we extract 267 biomarkers organized in six categories:

\begin{table}[H]
\centering
\small
\caption{Summary of 267 handcrafted EEG features.}
\label{tab:features}
\begin{tabular}{@{}lrc@{}}
\toprule
\textbf{Category} & \textbf{Dim} & \textbf{Description} \\
\midrule
Temporal statistics   & 57  & Mean, variance, IQR ($\times 19$ ch) \\
Spectral power (7 bands) & 133 & $\delta$, $\theta$, $\alpha_1$, $\alpha_2$, $\beta_1$, $\beta_2$, $\gamma$ \\
Entropy               & 38  & Permutation \& sample entropy \\
Connectivity           & 22  & Clustering, path length, efficiency \\
Asymmetry              & 3   & Interhemispheric differences \\
Kalman-filtered        & 14  & Smoothed PSD means \& residuals \\
\bottomrule
\end{tabular}
\end{table}

Spectral features are computed via Welch's method (PSD). Entropy is estimated using permutation entropy (order 3, normalized) and sample entropy (order 2) from the \texttt{antropy} library. Connectivity metrics include clustering coefficient, characteristic path length, global efficiency, and small-worldness derived from the Pearson correlation matrix. Kalman filtering smooths band power estimates to reduce non-stationarity.

\subsection{ADFormerHybrid Architecture}

The model consists of two parallel branches merged via cross-modal fusion (Figure~\ref{fig:arch}):

\subsubsection{Branch 1: Patch Transformer Encoder}

Raw EEG segments ($512 \times 19$) are divided into $P = 8$ temporal patches, each of dimension $64 \times 19 = 1{,}216$. A linear projection maps each patch to $d_{\text{model}} = 256$. Learnable positional embeddings are added:

\begin{equation}
  \mathbf{z}_i = W_{\text{proj}} \cdot \text{patch}_i + \mathbf{e}_i^{\text{pos}}, \quad i \in [1, P]
\end{equation}

The sequence is processed by a 4-layer Transformer encoder with 4 attention heads.

\subsubsection{Branch 2: Feature Encoder}

The 267-dimensional feature vector passes through:
\begin{equation}
  \mathbf{f} = W_2 \cdot \text{ReLU}(W_1 \cdot \text{LayerNorm}(\mathbf{x}_{\text{feat}})) + \text{dropout}
\end{equation}

yielding a $d_{\text{model}}$-dimensional representation.

\subsubsection{Cross-Modal Fusion}

The feature representation $\mathbf{f}$ is expanded and added to each patch token before the Transformer layers:
\begin{equation}
  \mathbf{z}_i' = \mathbf{z}_i + \mathbf{f}, \quad \forall i \in [1, P]
\end{equation}

This allows the self-attention mechanism to attend jointly to both raw temporal patterns and domain-specific biomarkers.

\subsubsection{Classification Head}

After Transformer encoding, tokens are mean-pooled and concatenated with $\mathbf{f}$:
\begin{equation}
  \hat{y} = \text{Head}\left(\left[\text{MeanPool}(\mathbf{Z}') \; \| \; \mathbf{f}\right]\right)
\end{equation}

where Head is: LayerNorm $\to$ Linear(512, 256) $\to$ ReLU $\to$ Dropout(0.1) $\to$ Linear(256, $C$), with $C \in \{2, 3\}$ depending on the task.

\subsection{Two-Stage Pipeline}

\begin{enumerate}[nosep]
  \item \textbf{Screening (binary):} AD vs.\ Healthy (2 classes)
  \item \textbf{Staging (ternary):} Mild vs.\ Moderate vs.\ Severe (3 classes), executed only if Stage 1 is positive
\end{enumerate}

Each stage uses an independent ADFormerHybrid model with task-specific weights.

\subsection{Ensemble Voting}

For each stage, five independently trained checkpoints (v1--v5) produce softmax probability vectors. The final prediction uses \textit{soft voting} (probability averaging):

\begin{equation}
  \hat{y}_{\text{ensemble}} = \arg\max_c \frac{1}{K} \sum_{k=1}^{K} p_k(c)
\end{equation}

where $K = 5$ and $p_k(c)$ is the softmax probability for class $c$ from model $k$.

% ═══════════ 4. Results ═══════════
\section{Experimental Results}

\subsection{Dataset and Training}

Models are trained on a clinical dataset of 65 patients with standard 19-channel EEG. Training uses cross-entropy loss with AdamW optimizer, learning rate $10^{-4}$, batch size 32, and standard scaling (per-feature z-normalization). Data augmentation includes temporal jittering and Gaussian noise injection.

\subsection{Performance}

\begin{table}[H]
\centering
\small
\caption{Performance comparison: single model vs.\ ensemble.}
\label{tab:results}
\begin{tabular}{@{}lccc@{}}
\toprule
\textbf{Configuration} & \textbf{Screen.\ Acc.} & \textbf{Stage Acc.} & \textbf{$\Delta$} \\
\midrule
Single model (best v.)  & 99.90\% & 96.64\% & --- \\
Ensemble (5 models)     & \textbf{99.95\%} & \textbf{97.79\%} & +1.15\% \\
\bottomrule
\end{tabular}
\end{table}

The ensemble provides a consistent improvement, particularly in staging where inter-class boundaries are more subtle. The screening task approaches ceiling performance, suggesting robust AD biomarker detection in the EEG signal.

\subsection{Inference Performance}

On Apple M1 (CPU-only), the full pipeline processes a 23~MB EEGLAB file (.set) containing 149 segments in approximately 7 seconds, including I/O, preprocessing, feature extraction, and dual-stage inference.

\subsection{Feature Importance}

SHAP analysis reveals that theta-band power, alpha power reduction, permutation entropy, and inter-hemispheric asymmetry are the most discriminative features, consistent with known AD neurophysiological markers \cite{jeong2004, dauwels2010}.

% ═══════════ 5. Platform ═══════════
\section{Neuro-Link Platform}

Beyond the model, we release \textbf{Neuro-Link}, a complete clinical decision support system built around ADFormerHybrid:

\begin{itemize}[nosep]
  \item \textbf{Backend:} FastAPI with async processing, security middleware (rate limiting, authentication, anti-injection), monitoring, PDF report export
  \item \textbf{Frontend:} React 18 with real-time status HUD, interactive radar visualizations, AI-generated clinical reports
  \item \textbf{AI Reports:} Gemini API integration with multi-model fallback chain for generating professional, accessible reports in French
  \item \textbf{Chatbot:} Conversational AI assistant to guide users through the analysis process
  \item \textbf{Deployment:} Docker, systemd, Nginx reverse proxy with TLS, CI/CD via GitHub Actions
\end{itemize}

The platform supports standard EEG formats: EEGLAB (.set), European Data Format (.edf), BrainVision (.vhdr), and OpenBCI CSV exports.

% ═══════════ 6. Limitations ═══════════
\section{Limitations and Future Work}

\begin{itemize}[nosep]
  \item \textbf{Dataset size:} 65 patients is limited for generalization claims. A multicenter study ($n > 500$) is planned (Phase 3).
  \item \textbf{No external validation:} Results are from internal evaluation. Independent replication is needed.
  \item \textbf{Regulatory status:} Not CE/FDA certified. Neuro-Link is a research tool, not a medical device.
  \item \textbf{Class imbalance:} The staging task may suffer from class imbalance in severity distribution.
\end{itemize}

Future work includes: multicenter clinical validation (CHU partnerships), CE MDR Class IIa certification, HL7/FHIR integration for hospital information systems, and longitudinal tracking of disease progression.

% ═══════════ 7. Conclusion ═══════════
\section{Conclusion}

We presented ADFormerHybrid, a dual-branch Transformer that fuses raw EEG patches with 267 handcrafted biomarkers through cross-modal token injection. Combined with soft-voting ensemble over five checkpoints, the model achieves 99.95\% screening and 97.79\% staging accuracy on clinical EEG data. The open-source Neuro-Link platform provides an accessible, secure, and deployable interface for EEG-based Alzheimer's screening research.

% ═══════════ Ethics ═══════════
\section*{Ethical Statement}

This work uses de-identified EEG data. Neuro-Link is a research tool and does not constitute a medical diagnostic device. All results must be interpreted by qualified healthcare professionals. The platform includes mandatory medical disclaimers.

% ═══════════ Code Availability ═══════════
\section*{Code and Data Availability}

Source code: \url{https://github.com/romainsantoli-web/neuro-link} (AGPL v3 license).\\
Model weights are available upon request due to file size constraints.

% ═══════════ References ═══════════
\begin{thebibliography}{10}

\bibitem{who2023}
World Health Organization, ``Dementia Fact Sheet,'' 2023. [Online]. Available: \url{https://www.who.int/news-room/fact-sheets/detail/dementia}

\bibitem{dauwels2010}
J.~Dauwels, F.~Vialatte, and A.~Cichocki, ``Diagnosis of Alzheimer's disease from EEG signals: Where are we standing?'' \textit{Current Alzheimer Research}, vol.~7, no.~6, pp.~487--505, 2010.

\bibitem{jeong2004}
J.~Jeong, ``EEG dynamics in patients with Alzheimer's disease,'' \textit{Clinical Neurophysiology}, vol.~115, no.~7, pp.~1490--1505, 2004.

\bibitem{ieracitano2020}
C.~Ieracitano, N.~Mammone, A.~Bramanti, A.~Hussain, and F.~C.~Morabito, ``A convolutional neural network approach for classification of dementia stages based on 2D-spectral representation of EEG recordings,'' \textit{Neurocomputing}, vol.~323, pp.~96--107, 2019.

\bibitem{morabito2017}
F.~C.~Morabito et al., ``Deep learning representation from electroencephalography of early-stage Creutzfeldt-Jakob disease and features for differentiation from rapidly progressive dementia,'' \textit{International Journal of Neural Systems}, vol.~27, no.~2, 2017.

\bibitem{dosovitskiy2020}
A.~Dosovitskiy et al., ``An image is worth 16x16 words: Transformers for image recognition at scale,'' in \textit{ICLR}, 2021.

\bibitem{lawhern2018}
V.~J.~Lawhern et al., ``EEGNet: A compact convolutional neural network for EEG-based brain--computer interfaces,'' \textit{Journal of Neural Engineering}, vol.~15, no.~5, 2018.

\bibitem{li2022}
Y.~Li et al., ``A hybrid deep learning framework for EEG-based Alzheimer's disease detection,'' \textit{IEEE Transactions on Neural Systems and Rehabilitation Engineering}, vol.~30, pp.~2570--2580, 2022.

\end{thebibliography}

\end{document}
